\begin{abstract}
Supervisory monitors in cyber-physical systems (CPS) need a timely and
logically consistent view of controller intent. On multicore platforms,
cache-coherent shared memory is an attractive export path, but
coherence does not make a multi-field control record atomic. This paper
studies publication over coherent memory as an architectural primitive
for CPS monitoring. Using gem5 ARM SE mode with Ruby
MESI\_Two\_Level, we compare five publication architectures:
\unsync, \seqlock, \dblbuf, \dmanaive, and \dmaseqlock. Across a
30-cell matrix with two PWM-inspired workloads and three contention
levels, the architectures without an explicit publication primitive
produce torn snapshots in every condition, whereas the disciplined
architectures do not. Among the correct designs, generation-based double
buffering over direct coherent sharing accepts 99.95\% of read attempts,
compared with 8.03\% for a synchronized DMA mirror, reduces average
CPU-side coherence traffic by 27\% relative to \seqlock, and shortens
whole-workload simulated time by 28\% relative to \dmaseqlock while
keeping retries negligible. Changing workload shape shifts absolute cost
but does not change the ranking: on the sustained-bias workload, the
direct-sharing protocols pay 35--37\% more CPU-side traffic, whereas
the synchronized DMA mirror stays nearly flat. The main conclusion is
that the decisive design variable is the publication primitive, not
transport alone.
\end{abstract}
